\subsection{Complete analytic closure of the $(8,3,7)$ budget}

\begin{theorem}[Proved: the jointly oriented $(8,3,7)$ piece]
\label{CP2:A-oriented-eighteen-eight-three-seven}
The budget $(|M|,|N|,|S|)=(8,3,7)$ contributes no
$\mathbb CP^2$ category in the joint orientation of
Theorem~\ref{CP2:A-two-object-eighteen-oriented-twenty-two}.
This conclusion uses a complete analytic partition of the source
catalog; no unfinished action enumeration is a premise.
\end{theorem}
\begin{proof}
The nonzero source restriction in the joint orientation requires
$H^2(BM;\mathbb Q)\ne0$.  The independently complete catalog of
such order-eight sources has $510$ classes.  Their proper minimal
ideals have the following shapes:
\[
\begin{array}{c|rrr}
\text{minimal band}&2\times2&2\times3&3\times2\\\hline
\text{source classes}&492&9&9.
\end{array}
\]
A nontrivial minimal local group is impossible here: two rows,
two columns and a local group of order at least two would already
give eight ideal elements and require an additional identity.
Among the seven ordinary monoids of order three, the previously
proved small-corner reductions leave exactly the identity-adjoined
two-element right-zero band $R_3$ as the opposite corner.

Every one of the $492$ two-by-two sources is covered by
Theorem~\ref{CP2:A-two-by-two-band-right-zero-all-states}.
In the literal catalog their actor images are $P_5$ or its
simultaneous row-and-column swap extension.  The actual quotient
identifies the minimal band bijectively, and the induced collage
map is an $H^2$ isomorphism.  All products therefore vanish.

All nine three-by-two sources have exactly two outside elements,
including the identity.  Theorem~\ref{CP2:A-two-column-two-outside-graph}
applies to them with the actual bounded homology of a hypothetical
$\mathbb CP^2$ collage.  It excludes every such source, for any
number of states.

Six of the two-by-three sources, with catalog indices
\[
492,\ 494,\ 495,\ 497,\ 499,\ 500,
\]
are covered by the all-state product theorem
\ref{CP2:A-eight-two-by-three-all-sources}.
The three remaining tables have literal arrow-reducing witnesses.
For source $493$, the elements $(a,e)=(7,1)$ are idempotent,
$ae=ea=e$, the sole-pair identification $a\sim e$ is a monoid
congruence, and $a-e$ is central in the integral monoid ring.
The proved integral idempotent-pair reduction gives a smaller
actual $\mathbb CP^2$ model from any hypothetical model here.
For sources $496$ and $498$, the idempotent $e=7$ satisfies
$Me=eMe$.  The extreme-idempotent corner reduction likewise
gives an actual homotopy-equivalent category with fewer arrows.
Each case contradicts the already proved unrestricted
seventeen-arrow lower bound.

These cases form the disjoint complete partition
\[
510=492+9+6+1+2.
\]
On the frozen refined complement of $69$ source tables, they
give precisely $56$ two-by-two cuts, seven two-column cuts,
and six explicit two-by-three cuts.  Thus that complement is
also exhausted exactly, independently of its unfinished
finite action stream.
\end{proof}

\begin{remark*}
The separate join
\path{monoid/certificates/eight_three_seven_analytic_complete.json}
binds the full literal source catalog to all proof receipts.
It checks $261\,120$ associativity equations and $31\,488$
actor-quotient multiplication equations, verifies the three
arrow-reduction witnesses, and compares the entire $510$-element
canonical set with the independently complete positive-$H^2$
catalog.  Its canonicalization enumerates every permutation of
minimal rows, minimal columns, and outside elements while fixing
the identity.  These are all possible isomorphisms because the
minimal ideal is characteristic.  The join also reconstructs
the unique required opposite type from all seven ordinary
order-three tables.  The earlier nineteen-budget snapshot is
preserved; this theorem supplies the twentieth completed piece,
leaving $(9,3,6)$ and $(10,3,5)$ pending.
\end{remark*}
